\chapter{Community weights}\label{ch:community}\index{community weights}

A community pack is a set of weights someone else published on Hugging Face that TandemLLM can load: a smaller quantisation, or a fine-tune of a model we serve. Since 0.8.2 the setup guide lists these packs under their model, after our own profiles, and \code{ops/use-profile.sh} runs one like any other profile. Listing a pack takes one JSON file in the repository. The first entry is ThinkingCap Qwen3.8-27B Kearuga by 0xWhiteMage, an EXL3 pack at 4.0 bits a weight.

\section{The packs in the setup guide}\label{sec:community-guide}

The weights step of the setup guide (\texttt{\#/setup}) shows a Community weights section below our profiles, one card per pack. A card shows the maker's GitHub avatar and \code{@login}, the networks the entry lists, the fine-tune it comes from, the format, the bits a weight and the download size. It also lists the modalities the pack holds, in grey with the reason when TandemLLM cannot serve one yet, and which of our drafters it runs with. Once a maintainer has run the pack, its card shows the measured speed too. A badge gives the license, another the status, and a link leads to the Hugging Face repository.

You can search over names, authors, handles, tags and the fine-tune's source, and sort by five keys, newest first by default. \enquote{Use} selects a pack the way a radio card selects one of our profiles, as \code{community/<id>}. Every later step takes that name, Wren included. A pack that cannot run images turns the Vision tool off, as the server would. \enquote{Add yours} opens the issue form of Section~\ref{sec:community-submit} with the model filled in.

\section{Where the list comes from}\label{sec:community-sync}\index{community weights!sync}

The list is \code{models/community-index.json}, one entry per file under \code{models/<model>/community/}. A box asks GitHub for that file on the public repository's \code{main} branch at two moments only: when the weights step opens and when you press Refresh (\code{POST /v1/setup/community/sync}). Nothing fetches it on a timer. The box sends the ETag of its last copy, so an unchanged list costs one 304 answer.

When the fetch fails, the page shows the box's last copy (\code{~/.tandemllm/community/index.json}). A box that never fetched one shows the list that came with its release. Look at the section's head to see which source the list came from and which entries it dropped. A 404 from GitHub is not an error: the repository has published no list yet, and the page shows a quiet note, \enquote{no published list yet}, beside the bundled list. A real failure is one short line without the URL, such as \enquote{GitHub could not be reached}; the server's log has the URL and the detail. \code{QSE_COMMUNITY_URL} points the box at another index, and \code{QSE_COMMUNITY_URL=off} never fetches.

A fetched index is untrusted. One with the wrong shape never replaces the cached copy, and an entry reaches the page only after it passes the registry's checks, those of Section~\ref{sec:community-verify}.

\section{Running a pack}\label{sec:community-run}

\begin{shell}
bash ~/TandemLLM/run/ops/use-profile.sh community/thinkingcap-27b-kearuga-exl3-4.00 --dry-run
bash ~/TandemLLM/run/ops/use-profile.sh community/thinkingcap-27b-kearuga-exl3-4.00
\end{shell}

The script writes the pack's profile from its entry to \code{run/community/<id>.env}, which \code{serve.env} reads. It downloads the pack at the commit the entry pins into \code{~/.tandemllm/community/<id>}, checks every weight file against the sha256 Hugging Face publishes, and fetches our draft models when they are missing. Then it stops the engine, starts it on the pack and sends one test chat. The checks are those of our own profiles: it refuses a download that would leave less than 60~GB of disk free, and it rolls back to the previous profile when the start or the test chat fails. \code{--dry-run} prints every URL and size and writes the profile into a folder of its own.

The memory estimate for a pack comes from our own profile of the same format and the nearest size: \code{exl3-4.0} for a 4.0 bpw EXL3 pack. It stays an estimate until a maintainer has loaded the pack. Since 0.12.1 a pack can be built on Nemotron 3.5 Super VL too (Section~\ref{sec:community-supervl}).

\begin{table}[htbp]
\centering\small
\caption[The first community pack]{The first community pack. Speed: one run per kind of text on a DGX Spark with 0.8.2, as the entry's \code{verified} block records it (\measured).}\label{tab:community-first}
\begin{tabularx}{\linewidth}{@{}L{.3\linewidth}Y@{}}
\toprule
\hd{field} & \hd{value} \\
\midrule
id & \code{thinkingcap-27b-kearuga-exl3-4.00} \\
author & 0xWhiteMage \\
fine-tune of & \code{bottlecapai/ThinkingCap-Qwen3.8-27B} \\
pack & \code{0xWhiteMage/ThinkingCap-Qwen3.8-27B-Kearuga-EXL3-4.00bpw}, EXL3 at 4.0 bits a weight, 16.38~GB, standalone \\
modalities & text; the pack also holds a 6-bit vision tower, which the standalone path does not serve yet \\
drafters & ours, the EXL3 K3 drafters and K4 draft head \\
speed & 39.6~tok/s prose, 95.9 code, 51.3 chat; 15.2 without speculation \\
license & PolyForm-Small-Business-1.0.0, with Apache-2.0 for Qwen's part \\
\bottomrule
\end{tabularx}
\end{table}

\section{Packs on Nemotron 3.5 Super VL}\label{sec:community-supervl}\index{community weights!Super VL}

A pack on Super VL loads on NemotronH (Chapter~\ref{ch:nemotron}) and is served with our EXL3 profile nearest in size: \code{nemotron-super-vl-exl3-4.0} for a 3.5 bpw pack. It takes that profile's memory limits, its sampler, NVIDIA's LiLiCorr drafter and the lookup, with the pack in place of the weights; the entry's \code{drafters} is \code{none}. The pack must be standalone, a NemotronH checkpoint with 88 layers, in EXL3. \code{python -m server.setup.community check-pack community/<id>} checks a downloaded pack's files and names what differs: the config, the layer count, the quantisation, every shard the index names, and the image tower when the entry lists images.

\begin{table}[htbp]
\centering\small
\caption[The first pack on Super VL]{The first pack on Super VL. Speed: one run per kind of text on a DGX Spark with 0.12.1, 512 greedy tokens, thinking off, as the entry's \code{verified} block records it (\measured); \code{models/nemotron-3.5-super-vl-120b-a12b/measurements/0.12.1-summary.json} has the runs.}\label{tab:community-supervl}
\begin{tabularx}{\linewidth}{@{}L{.3\linewidth}Y@{}}
\toprule
\hd{field} & \hd{value} \\
\midrule
id & \code{nemotron-super-vl-exl3-3.5bpw} \\
author & styles01 (pull request \#63) \\
pack & \code{aitamedia/NVIDIA-Nemotron-3.5-Super-VL-120B-A12B-EXL3-3.5bpw}, EXL3: routed experts at 3.5 bits, every other projection at 5, an fp16 head; 57.8~GB, standalone \\
modalities & text; the pack holds no image tower \\
drafters & LiLiCorr and the lookup, from \code{nemotron-super-vl-exl3-4.0} \\
speed & 33.2~tok/s prose, 53.9 code, 37.4 chat, 99.0 edit, 112.3 copy; 25.9 without speculation; every drafted text the plain run's \\
agreement & top token equal to exllamav3 1.6.0's own forward on the pack at 510 of 512 positions, both others near ties; KL 0.0019 \\
memory & 53.58~GiB of weights at load; the served profile (262,144 tokens, 3 slots) starts and answers in 110~s \\
license & OpenMDW-1.1 \\
\bottomrule
\end{tabularx}
\end{table}

Plain decoding runs at 25.9~tok/s on this pack and at 31.1 on our 4.0 pack (0.11.0). From the headers, a decode token reads 7.95~GB of this pack and 6.74~GB of ours.

\section{Submitting a pack}\label{sec:community-submit}

Send a pull request that adds one file: \code{models/<model>/community/<id>.json}. The file name is the id. The first entry, \code{models/qwen3.8-27b/community/thinkingcap-27b-kearuga-exl3-4.00.json}, is a good one to copy, and the JSON Schema is \code{models/schema/community.schema.json}. If you would rather not write JSON, open an issue with the \enquote{Add a community weight pack} form (\code{.github/ISSUE_TEMPLATE/community-quant.yml}) and a maintainer writes the file.

The fields that matter most:

\begin{itemize}
\item \code{hf.repo} and \code{hf.revision}: a public repository, and the full 40-character commit. A branch name is refused, because a branch can change after the review.
\item \code{format} and \code{bpw}: only \code{exl3} passes today. The schema knows three more formats, which the check refuses until TandemLLM can load them.
\item \code{load}: \code{standalone} when the pack is the whole model (Section~\ref{sec:exl3-standalone}), or \code{overlay} when it holds EXL3 projections that go over the base checkpoint.
\item \code{modalities}: what works in TandemLLM today. A standalone pack is text only for now, even with a vision tower inside; list \code{image} in \code{modalities_in_pack} instead. An overlay keeps the base checkpoint's vision tower.
\item \code{drafters}: \code{tandem-exl3} for our EXL3 drafters; \code{docs/community-weights.md} lists the other two values.
\item \code{author.github}, your GitHub user name, whose avatar the card shows. Leave \code{status} at \code{submitted}.
\end{itemize}

\code{docs/community-weights.md} lists every field. The registry refuses a pack that is private or gated, a branch in place of a commit, a format TandemLLM cannot load, a base model other than \code{qwen3.8-27b} and \code{nemotron-3.5-super-vl-120b-a12b}, a modality the load path cannot serve, a link that is not https, markup in a text field, and a license that does not let others download and run the weights.

\section{How verification works}\label{sec:community-verify}

Verification has two parts. On the pull request, the \code{community} check runs \code{python -m server.setup.registry --check} and \code{python -m server.setup.community --index --check}, which hold the file to the schema and to the rules above. Then a maintainer downloads the pack at the pinned commit and runs it on a DGX Spark. If it loads and answers, the maintainer fills \code{verified} (the build, the device, the date, tok/s per kind of text, tokens per round, a note) and sets \code{status} to \code{verified}. A pack that does not load stays \code{submitted}, or becomes \code{broken} with a note.

After the merge, a maintainer regenerates the index (\code{python -m server.setup.community --index}). Every box sees the pack the next time someone opens the weights step or presses Refresh.
